2022 考研数学基础班


\section*{第八章 曲线积分与曲面积分}

- 曲线积分
- 第一类曲线积分
- 第二类曲线积分 —— 两类曲线积分之间的联系
- 格林公式 —— 积分与路径无关的条件
- 曲面积分
- 第一类曲面积分
- 第二类曲面积分 —— 两类曲面积分之间的联系
- 高斯公式 —— 散度
- 斯托克斯公式 —— 旋度

第一类曲线积分的定义
设 L 为 xOy 面内的一条光滑曲线弧, 函数 f(x,y) 在 L 上有界. 在 L 上任意插入一点列 M₁, M₂, ..., Mₙ₋₁ 把 L 分成 n 个小段. 设第 i 个小段的长度为 Δsi. 又 (ξᵢ, ηᵢ) 为第 i 个小段上任意取定的一点, 作乘积 f(ξᵢ, ηᵢ)Δsi(i = 1, 2, ..., n), 并作和 ∑ⁿᵢ=₁ f(ξᵢ, ηᵢ)Δsi. 若当各小弧段的长度最大值 λ → 0 时, 这和极限总存在, 且与曲线弧 L 的分法及点 (ξᵢ, ηᵢ) 的取法无关, 则称此极限为函数 f(x,y) 在曲线弧 L 上对弧长的曲线积分或第一类曲线积分 (以后我们均采用第二种名称), 记作 ∫ₗ L f(x,y)ds, 即
\[ \int_L f(x,y) ds = \lim_{\lambda \to 0} \sum_{i=1}^n f(\xi_i, \eta_i) \Delta s_i, \]
其中 f(x,y) 叫做被积函数, L 叫做积分弧段.

第一类曲线积分的性质
(1) 线性性 设 \(\alpha, \beta\) 为常数, 则
\[ \int_L [\alpha f(x, y) + \beta g(x, y)] ds = \alpha \int_L f(x, y) ds + \beta \int_L g(x, y) ds. \]
(2) 可加性 若积分弧段 \(L\) 可划分成两段光滑曲线弧 \(L_1\) 和 \(L_2\), 则
\[ \int_L f(x, y) ds = \int_{L_1} f(x, y) ds + \int_{L_2} f(x, y) ds. \]
(3) 若在区域 \(D\) 上, \(f(x, y) \equiv 1\), \(L\) 为 \(D\) 内一条分段光滑曲线, \(l\) 为 \(L\) 的长度, 则
\[ l = \int_L ds. \]

(4) 设在 L 上, \(f(x,y) \leq g(x,y)\), 则
\[ \int_L f(x,y) \, \mathrm{d}s \leq \int_L g(x,y) \, \mathrm{d}s. \]
特别地, 有
\[ \left| \int_L f(x,y) \, \mathrm{d}s \right| \leq \int_L |f(x,y)| \, \mathrm{d}s. \]
第一类曲线积分的计算
基本方法 将积分弧段参数化后把曲线积分化为定积分.

设 \(f(x,y)\) 在曲线弧 \(L\) 上有定义且连续,
(1) 曲线 \(L\) 由参数方程表示的情形 \(L\) 的参数方程为 \(x = \varphi(t)\), \(y = \psi(t) (\alpha \leq t \leq \beta)\). 若 \(\varphi(t)\), \(\psi(t)\) 在 \([\alpha, \beta]\) 上具有一阶连续导数, 且 \([\varphi'(t)]^2 + [\psi'(t)]^2 \neq 0\), 则曲线积分 \(\int_L f(x,y) ds\) 存在, 且
\[ \int_L f(x,y) ds = \int_\alpha^\beta f(\varphi(t), \psi(t)) \sqrt{[\varphi'(t)]^2 + [\psi'(t)]^2} dt. \]

(2)曲线 L 由 \(y = \psi(x)(x_0 \le x \le x_1)\) 表示的情形 可以把这种情况看作特殊的参数方程: \(x = x, y = \psi(x)(x_0 \le x \le x_1)\), 从而有
\[ \int_L f(x, y) ds = \int_{x_0}^{x_1} f(x, \psi(x)) \sqrt{1 + [\psi'(x)]^2} dx. \]
(3)曲线 L 由 \(x = \varphi(y)(y_0 \le y \le y_1)\) 表示的情形与 (2) 类似.(见讲义)
(4)曲线 L 由极坐标形式 \(r = r(\theta)(\alpha \le \theta \le \beta)\) 表示的情形
\[ \int_L f(x, y) ds = \int_{\alpha}^{\beta} f(r(\theta) \cos \theta, r(\theta) \sin \theta) \sqrt{[r(\theta)]^2 + [r'(\theta)]^2} d\theta. \]
(5)空间曲线弧情形.(见讲义)

例1
设l为椭圆 \(\frac{x^2}{4} + \frac{y^2}{3} = 1\), 其周长为a, 则 \(\oint (2xy + 3x^2 + 4y^2)ds = \_\_\_\_\_\_\_\_\_\_\_\_\_\_\_\_\_\_\_\_\_.

例 2 求 \(\int_{L} y ds\), 其中 \(L\) 为抛物线 \(y^2 = x\) 上连接点 \(A(0, 0)\) 到点 \(B(1, -1)\) 的曲线弧.

第一类曲线积分的对称性
积分弧段关于坐标面对称.(见讲义)
积分弧段关于变量具有轮换对称性.
设分段光滑的空间曲线 Γ 对变量 x, y, z 具有轮换对称性 (即轮换坐标轴的名称, 曲线 Γ 的方程不变), f(x, y, z) 在 Γ 上连续, 则
\[
\begin{align*}
\int_{\Gamma} f(x, y, z) \, \mathrm{d}s &= \int_{\Gamma} f(y, z, x) \, \mathrm{d}s = \int_{\Gamma} f(z, x, y) \, \mathrm{d}s \\
&= \frac{1}{3} \int_{\Gamma} [f(x, y, z) + f(y, z, x) + f(z, x, y)] \, \mathrm{d}s.
\end{align*}
\]

例3 设L为球面x²+y²+z²=1与平面x+y+z=0的交线,则ΦLz²ds=

第二类曲线积分的定义
设 L 为 xOy 面内从点 A 到点 B 的一条有向光滑曲线弧, 函数 P(x, y) 与 Q(x, y) 在 L 上有界. 在 L 上沿 L 的方向任意插入一点列 M₁(x₁, y₁), M₂(x₂, y₂), · · · , Mₙ₋₁(xₙ₋₁, yₙ₋₁), 把 L 分成 n 个有向小弧段
\[ \tilde{M}_{i-1} \tilde{M}_i (i = 1, 2, \dots, n; M_0 = A, M_n = B). \]
设 Δxᵢ = xᵢ - xᵢ₋₁, Δyᵢ = yᵢ - yᵢ₋₁, 点 (ξᵢ, ηᵢ) 为 Mᵢ₋₁ Mᵢ 上任意取定的点, 作乘积 P(ξᵢ, ηᵢ) Δxᵢ (i = 1, 2, · · · , n), 并作和 ∑ⁿᵢ₋₁ P(ξᵢ, ηᵢ) Δxᵢ. 若当各小弧段长度的最大值 λ → 0 时, 这和极限总存在, 且与曲线弧 L 上的方法及点 (ξᵢ, ηᵢ) 的取法无关, 则称此极限为函数 P(x, y) 在有向曲线弧 L 上对坐标 x 的曲线积分, 记作 ∫ₗ P(x, y)dx.

第二类曲线积分的物理意义
在应用上, 第二类曲线积分经常以向量形式出现. 令向量值函数 \(\mathbf{F}(x, y) = P(x, y)\mathbf{i} + Q(x, y)\mathbf{j}\), 积分弧段为 \(L\). 注意到 \(\mathrm{d}\mathbf{r} = \mathrm{d}x\mathbf{i} + \mathrm{d}y\mathbf{j}\), 若将 \(\mathbf{F}\) 看作变力, \(L\) 看作 \(\mathbf{F}\) 做功的路径, 则
\[ \int_L \mathbf{F}(x, y) \cdot \mathrm{d}\mathbf{r} = \int_L P(x, y) \mathrm{d}x + Q(x, y) \mathrm{d}y \]
表示变力 \(\mathbf{F}\) 沿路径 \(L\) 所做的功.

第一类曲线积分与第二类曲线积分的比较
<table>曲线积分积分元素积分范围是否有向典型物理意义∫<sub>L</sub> f(x,y)ds弧长元素曲线弧不带方向曲线形构件的质量∫<sub>L</sub> F(x,y)·dr有向曲线元有向曲线弧带方向变力沿曲线所做的功</table>
第二类曲线积分的性质 (见讲义) 线性性、可加性、有向性.

第二类曲线积分的计算
设 \(P(x,y)\) 与 \(Q(x,y)\) 在有向曲线弧 \(L\) 上有定义且连续， \(L\) 的参数方程为 \(\left\{ \begin{array}{l}x = \phi (t),\\ y = \psi (t). \end{array} \right.\) 当参数 \(t\) 单调地由 \(\alpha\) 变到 \(\beta\) 时，点 \(M(x,y)\) 从 \(L\) 的起点 \(A\) 沿 \(L\) 运动到终点 \(B\) . 若 \(\phi (t)\) 与 \(\psi (t)\) 在以 \(\alpha\) 和 \(\beta\) 为端点的闭区间上具有一阶连续导数，且 \([\phi '(t)]^2 +[\psi '(t)]^2 \neq 0\) ，则曲线积分 \(\int_{L}P(x,y)\mathrm{d}x + Q(x,y)\mathrm{d}y\) 存在，且
\[\int_{L}P(x,y)\mathrm{d}x + Q(x,y)\mathrm{d}y = \int_{\alpha}^{\beta}[P(\phi (t),\psi (t))\phi '(t) + Q(\phi (t),\psi (t))\psi '(t)]\mathrm{d}t.\]
在计算时需要注意，下限 \(\alpha\) 对应 \(L\) 的起点，上限 \(\beta\) 对应 \(L\) 的终点， \(\alpha\) 不一定小于 \(\beta\) .

例 4 计算 \(\int_{L} xydx\)，其中 \(L\) 为抛物线 \(y^2 = x\) 上从点 \(A(1, -1)\) 到点 \(B(1, 1)\) 的一段弧.

两类曲线积分之间的联系
平面曲线弧 L 上的两类曲线积分之间有如下联系:
\[ \int_L Pdx + Qdy = \int_L (P \cos \alpha + Q \cos \beta) ds, \]
其中 \(\alpha(x, y)\), \(\beta(x, y)\) 为有向曲线弧 L 在点 \((x, y)\) 处的切向量的方向角.
类似地, 空间曲线弧 Γ 上的两类曲线积分之间有如下联系:
\[ \int_{\Gamma} Pdx + Qdy + Rdz = \int_{\Gamma} (P \cos \alpha + Q \cos \beta + R \cos \gamma) ds, \]
其中 \(\alpha(x, y, z)\), \(\beta(x, y, z)\), \(\gamma(x, y, z)\) 为有向曲线弧 Γ 在点 \((x, y, z)\) 处的切向量的方向角.

例5 把对坐标的曲线积分 \(\int_{L} P(x, y) dx + Q(x, y) dy\) 化成对弧长的曲线积分, 其中 \(L\) 为沿上半圆周 \(x^2 + y^2 = 2x\) 从点 \((0, 0)\) 到点 \((1, 1)\) 的曲线弧.

曲线积分的应用 (见讲义)
第一类曲线积分的应用.
第二类曲线积分的应用.
例 6
质点 P 沿着以 AB 为直径的半圆周, 从点 A(1,2) 运动到点 B(3,4) 的过程中受变力 F 作用 (如图), F 的大小等于点 P 与原点 O 之间的距离, 其方向垂直于线段 OP, 且与 y 轴正向的夹角小于 \(\frac{\pi}{2}\). 求变力 F 对质点 P 所做的功.

同步习题 见讲义第一节同步习题.

第二节 曲线积分（二）
## 格林公式
设闭区域 D 由分段光滑的曲线 L 围成, 若函数 P(x,y) 及 Q(x,y) 在 D 上具有一阶连续偏导数, 则有
\[ \iint_D \left( \frac{\partial Q}{\partial x} - \frac{\partial P}{\partial y} \right) dxdy = \oint_L Pdx + Qdy, \]
其中 L 是 D 的取正向的边界曲线.

例7 计算曲线积分 \(\int_{L} \sin 2x dx + 2(x^2 - 1) y dy\), 其中 \(L\) 是曲线 \(y = \sin x\) 上从点 \((0, 0)\) 到点 \((\pi, 0)\) 的一段.

例8 计算曲线积分 \(I = \oint_{L} \frac{xdy - ydx}{4x^2 + y^2}\)，其中 \(L\) 是以点 \((1,0)\) 为中心，\(R\) 为半径的圆周 \((R > 1)\)，取逆时针方向。

积分与路径无关的条件
设 \(D\) 为平面中的单连通区域,函数 \(P(x,y)\) 和 \(Q(x,y)\) 在 \(D\) 上具有一阶连续偏导数,则以下条件相互等价:
(1)对于 \(D\) 中任意分段光滑闭曲线 \(C\) ,都有
\[\oint_{C}P\mathrm{d}x + Q\mathrm{d}y = 0;\]
(2)对于 \(D\) 中从点 \(M_{1}\) 到点 \(M_{2}\) 的任意两条分段光滑曲线 \(L_{1}\) , \(L_{2}\) ,都有
\[\int_{L_{1}}P\mathrm{d}x + Q\mathrm{d}y = \int_{L_{2}}P\mathrm{d}x + Q\mathrm{d}y;\]
(3)存在在 \(D\) 上具有一阶连续偏导数的函数 \(u(x,y)\) ,使得 \(\mathrm{d}[u(x,y)] = P\mathrm{d}x + Q\mathrm{d}y\) , \(u(x,y)\) 被称为微分式 \(P\mathrm{d}x + Q\mathrm{d}y\) 在 \(D\) 中的一个原函数;
(4)在 \(D\) 中,恒有
\[\frac{\partial Q}{\partial x} = \frac{\partial P}{\partial y}.\]

例9 若曲线积分 \(\int_{L} \frac{xdx - aydy}{x^2 + y^2 - 1}\) 在区域 \(D = \{(x, y) | x^2 + y^2 < 1\}\) 内与路径无关, 则 \(a =\) ______.

曲线积分的基本定理
设 \(F(x,y) = P(x,y)i + Q(x,y)j\) 是平面区域 \(D\) 内的一个向量场. 若 \(P(x,y)\) 与 \(Q(x,y)\) 都在 \(D\) 内连续, 且存在一个数量函数 \(f(x,y)\), 使得
\[F = \nabla f = \frac{\partial f}{\partial x}i + \frac{\partial f}{\partial y}j,\]
则曲线积分 \(\int_L F \cdot dr\) 在 \(D\) 内与路径无关, 且
\[\int_L F \cdot dr = f(B) - f(A),\]
其中 \(L\) 是位于 \(D\) 内起点为 \(A\), 终点为 \(B\) 的任一分段光滑曲线.

例 10
设函数 \(f(x,y)\) 满足 \(\frac{\partial f(x,y)}{\partial x} = (2x+1)e^{2x-y}\), 且 \(f(0,y) = y+1\), \(L_t\) 是从点 \((0,0)\) 到点 \((1,t)\) 的光滑曲线. 计算曲线积分 \(I(t) = \int_{L_t} \frac{\partial f(x,y)}{\partial x} dx + \frac{\partial f(x,y)}{\partial y} dy\), 并求 \(I(t)\) 的最小值.

全微分的原函数
(1) \(P(x,y)dx + Q(x,y)dy\) 为全微分的充分必要条件 设区域 \(D\) 是一个单连通区域, 若函数 \(P(x,y)\) 与 \(Q(x,y)\) 在 \(D\) 内具有一阶连续偏导数, 则 \(P(x,y)dx + Q(x,y)dy\) 在 \(D\) 内为某一函数 \(u(x,y)\) 的全微分的充分必要条件是 \(\frac{\partial Q}{\partial x} = \frac{\partial P}{\partial y}\) 在 \(D\) 内恒成立.
(2) 全微分方程 若一个微分方程能写成
\[P(x,y)dx + Q(x,y)dy = 0\]
的形式, 而 \(P(x,y)dx + Q(x,y)dy\) 为某一个函数 \(u(x,y)\) 的全微分, 则上述方程称为全微分方程, \(u(x,y) = C\) 是它的隐式通解, 其中 \(C\) 为任意常数.
已知某微分式, 我们求其原函数 \(u(x,y)\) 的方法一般有如下三种:
(1) 折线法; (2) 积分法; (3) 凑微分法.

例 11
验证 \((\frac{y}{x} + x^2) dx + (\ln x - 2y) dy\) 在右半平面 \(x > 0\) 内是某一函数 \(u(x, y)\) 的全微分, 并求这样的一个 \(u(x, y)\).

同步习题 见讲义第二节同步习题.

第一类曲面积分的定义
设曲面 \(\Sigma\) 是光滑的,函数 \(f(x,y,z)\) 在 \(\Sigma\) 上有界.把 \(\Sigma\) 任意分成 \(n\) 小块 \(\Delta S_{i}(\Delta S_{i}\) 同时也代表第 \(i\) 块小曲面的面积),设 \((\xi_{i},\eta_{i},\zeta_{i})\) 是 \(\Delta S_{i}\) 上任意取定的一点,作乘积 \(f(\xi_{i},\eta_{i},\zeta_{i})\Delta S_{i}(i = 1,2,\dots ,n)\) ,并作和 \(\sum_{i = 1}^{n}f(\xi_{i},\eta_{i},\zeta_{i})\Delta S_{i}\) 若当各小块曲面的直径的最大值 \(\lambda \rightarrow 0\) 时,这和极限总存在,且与曲面 \(\Sigma\) 的 分法及点 \((\xi_{i},\eta_{i},\zeta_{i})\) 的取法无关,则称此极限为函数 \(f(x,y,z)\) 在曲面 \(\Sigma\) 上对 面积的曲面积分或第一类曲面积分,记作 \(\iint_{\Sigma}f(x,y,z)\mathrm{d}S\) ,即
\[\iint_{\Sigma}f(x,y,z)\mathrm{d}S = \lim_{\lambda \to 0}\sum_{i = 1}^{n}f(\xi_{i},\eta_{i},\zeta_{i})\Delta S_{i},\]
其中 \(f(x,y,z)\) 叫做被积函数, \(\Sigma\) 叫做积分曲面.

第一类曲面积分的计算
若积分曲面 Σ 由方程 z = z(x, y) 给出, Σ 在 xOy 面上的投影区域为 Dxy, 函数 z = z(x, y) 在 Dxy 上具有一阶连续偏导数, 被积函数 f(x, y, z) 在 Σ 上连续, 则
\[ \iint_{\Sigma} f(x, y, z) \mathrm{d}S = \iint_{D_{xy}} f(x, y, z(x, y)) \sqrt{1 + [z_x'(x, y)]^2 + [z_y'(x, y)]^2} \mathrm{d}x \mathrm{d}y. \]
第一类曲面积分的对称性 (见讲义)

例 12 求抛物面壳 \(z = \frac{1}{2}(x^2 + y^2)(0 \le z \le 1)\) 的质量, 此壳的面密度为 \(\mu = z\).

第二类曲面积分的定义
设 Σ 为光滑的有向曲面, 函数 R(x, y, z) 在 Σ 上有界. 把 Σ 任意分成 n 块小曲面 ΔSᵢ(ΔSᵢ 同时也代表第 i 块小曲面的面积), ΔSᵢ 在 xOy 面上的投影为 (ΔSᵢ)ₓᵧ, (ξᵢ, ηᵢ, ζᵢ) 为 ΔSᵢ 上任意取定的一点, 作乘积 R(ξᵢ, ηᵢ, ζᵢ)(ΔSᵢ)ₓᵧ(i = 1, 2, ..., n), 并作和
∑ n R(ξi, ηi, ζi)(ΔSi)xy. 若当各小块曲面的直径的最大值 λ → 0 时, 这和的极限总存在, 且与曲面 Σ 的分法以及点 (ξi, ηi, ζi) 的取法无关, 则称此极限为函数 R(x, y, z) 在有向曲面 Σ 上对坐标 x, y 的曲面积分, 记作 ∬ Σ R(x, y, z)dxdy, 即
\[ \iint_{\Sigma} R(x, y, z) dxdy = \lim_{\lambda \to 0} \sum_{i=1}^{n} R(\xi_i, \eta_i, \zeta_i)(\Delta S_i)_{xy}, \]
其中 R(x, y, z) 叫做被积函数, Σ 叫做积分曲面.

右手直角坐标系下有向曲面、法向量以及方向余弦的关系
<table>曲面的表示方程曲面的侧法向量方向余弦<td rowspan="2">z = z(x, y)上侧(-z'<sub>x</sub>, -z'<sub>y</sub>, 1)cos γ &gt; 0下侧(z'<sub>x</sub>, z'<sub>y</sub>, -1)cos γ &lt; 0<td rowspan="2">y = y(z, x)右侧(-y'<sub>x</sub>, 1, -y'<sub>z</sub>)cos β &gt; 0左侧(y'<sub>x</sub>, -1, y'<sub>z</sub>)cos β &lt; 0<td rowspan="2">x = x(y, z)前侧(1, -x'<sub>y</sub>, -x'<sub>z</sub>)cos α &gt; 0后侧(-1, x'<sub>y</sub>, x'<sub>z</sub>)cos α &lt; 0</table>

第二类曲面积分的计算
设积分曲面 \(\Sigma\) 由方程 \(z = z(x,y)\) 给出， \(\Sigma\) 在 \(xOy\) 面上的投影区域为 \(D_{xy}\) ，函数 \(z = z(x,y)\) 在 \(D_{xy}\) 上具有一阶连续偏导数，被积函数 \(R(x,y,z)\) 在 \(\Sigma\) 上连续.于是
\[\iint_{\Sigma}R(x,y,z)\mathrm{d}x\mathrm{d}y = \pm \iint_{D_{xy}}R(x,y,z(x,y))\mathrm{d}x\mathrm{d}y.\]
若 \(\Sigma\) 取上侧，即 \(\cos \gamma >0\) ，则上式右端取正号；若 \(\Sigma\) 取下侧，即 \(\cos \gamma < 0\) ，则上式右端取负号.
第二类曲面积分为零的三种特殊情况：
当曲面 \(\Sigma\) 垂直于 \(xOy\) 面时， \(\iint_{\Sigma}R(x,y,z)\mathrm{d}x\mathrm{d}y = 0\) 当曲面 \(\Sigma\) 垂直于 \(yOz\) 面时， \(\iint_{\Sigma}P(x,y,z)\mathrm{d}y\mathrm{d}z = 0\) 当曲面 \(\Sigma\) 垂直于 \(zOx\) 面时， \(\iint_{\Sigma}Q(x,y,z)\mathrm{d}z\mathrm{d}x = 0\)

两类曲面积分之间的联系 两类曲面积分之间的联系如下：
\[ \iint_{\Sigma} Pdy dz + Qdz dx + Rdx dy = \iint_{\Sigma} (P \cos \alpha + Q \cos \beta + R \cos \gamma) dS, \]
其中 \(\cos \alpha\), \(\cos \beta\), \(\cos \gamma\) 为有向曲面 \(\Sigma\) 在点 \((x, y, z)\) 处的法向量的方向余弦.

例 13
计算曲面积分 \(\iint_S \frac{x dy dz + z^2 dx dy}{x^2 + y^2 + z^2}\), 其中 \(S\) 是由曲面 \(x^2 + y^2 = R^2\) 及两平面 \(z = R\), \(z = -R(R > 0)\) 所围成立体表面的外侧.

例 14
计算
∬[f(x,y,z)+x]dydz+[2f(x,y,z)+y]dzdx+[f(x,y,z)+z]dxdy, 其中 f(x,y,z) 为连续函数, Σ 是平面 x-y+z=1 在第四卦限部分的上侧.

同步习题 见讲义第三节同步习题.

第四节 曲面积分（二）
## 高斯公式
设空间闭区域 Ω 由分片光滑的闭曲面 Σ 所围成. 若函数 P(x, y, z), Q(x, y, z), R(x, y, z) 在 Ω 上具有一阶连续偏导数, 则有
\[ \iiint_{\Omega} \left( \frac{\partial P}{\partial x} + \frac{\partial Q}{\partial y} + \frac{\partial R}{\partial z} \right) dv = \iint_{\Sigma} P dy dz + Q dz dx + R dx dy, \]
或
\[ \iiint_{\Omega} \left( \frac{\partial P}{\partial x} + \frac{\partial Q}{\partial y} + \frac{\partial R}{\partial z} \right) dv = \iint_{\Sigma} (P \cos \alpha + Q \cos \beta + R \cos \gamma) dS. \]
这里 Σ 是 Ω 的整个边界曲面的外侧, cos α, cos β, cos γ 是 Σ 在点 (x, y, z) 处的法向量的方向余弦.

例 15
计算 \(\iint_{\Sigma} \frac{axdy dz + (z+a)^2 dxdy}{(x^2 + y^2 + z^2)^{\frac{1}{2}}}\)，其中 \(\Sigma\) 为下半球面 \(z = -\sqrt{a^2 - x^2 - y^2}\) 的上侧，\(a\) 为大于零的常数。

例 16
计算曲面积分 I = ∫Σ x d y d z + y d z d x + z d x d y , 其中 Σ 是曲面 2 x 2 + 2 y 2 + z 2 = 4 的外侧.

散度
设有向量场
\[A(x,y,z) = P(x,y,z)i + Q(x,y,z)j + R(x,y,z)k,\]
其中函数 \(P\) , \(Q\) , \(R\) 均具有一阶连续偏导数, \(\Sigma\) 是场内一片有向曲面, \(n\) 是 \(\Sigma\) 在点 \((x,y,z)\) 处的单位法向量, 则积分
\[\iint_{S}A\cdot ndS\]
称为向量场 \(A\) 通过曲面 \(\Sigma\) 向着指定侧的通量 (或流量).
\(\frac{\partial P}{\partial x} +\frac{\partial Q}{\partial y} +\frac{\partial R}{\partial z}\) 叫做向量场 \(A\) 的散度, 记作 div \(A\) . 利用向量微分算子 \(\nabla = \left(\frac{\partial}{\partial x}, \frac{\partial}{\partial y}, \frac{\partial}{\partial z}\right)\) , \(A\) 的散度 div \(A\) 也可以表达为 \(\nabla \cdot A\) , 即
\[\operatorname {div} A = \nabla \cdot A.\]

例 17 向量场 \(u(x, y, z) = xy^2i + ye^z j + x \ln(1 + z^2)k\) 在点 \(P(1, 1, 0)\) 处的散度 \(\text{div } u = \_\_\_\_\_\_\_\_\_\_\_\_\_\_\_\_\_\_\_\_\_\).

斯托克斯公式
设 \(\Gamma\) 为分段光滑的空间有向闭曲线， \(\Sigma\) 为以 \(\Gamma\) 为边界的分片光滑的有向曲面， \(\Gamma\) 的正向与 \(\Sigma\) 的侧符合右手规则（当右手除拇指外的四指沿 \(\Gamma\) 的绕行方向时，拇指所指的方向与 \(\Sigma\) 上法向量的指向相同，此时称 \(\Gamma\) 是有向曲面 \(\Sigma\) 的正向边界曲线).若函数 \(P(x,y,z)\) ， \(Q(x,y,z)\) ， \(R(x,y,z)\) 在曲面 \(\Sigma\) (连同边界 \(\Gamma\) )上具有一阶连续偏导数，则有
\[\iint_{\Sigma}\left(\frac{\partial R}{\partial y} -\frac{\partial Q}{\partial z}\right)\mathrm{d}y\mathrm{d}z + \left(\frac{\partial P}{\partial z} -\frac{\partial R}{\partial x}\right)\mathrm{d}z\mathrm{d}x + \left(\frac{\partial Q}{\partial x} -\frac{\partial P}{\partial y}\right)\mathrm{d}x\mathrm{d}y\] \[= \oint_{\Gamma}P\mathrm{d}x + Q\mathrm{d}y + R\mathrm{d}z.\]

可以利用行列式记号把斯托克斯公式写成
\[ \iint_{\Sigma} \begin{vmatrix} \frac{\mathrm{d}y\mathrm{d}z}{\partial x} & \frac{\mathrm{d}z\mathrm{d}x}{\partial y} & \frac{\mathrm{d}x\mathrm{d}y}{\partial z} \\ \frac{\partial}{\partial x} & \frac{\partial}{\partial y} & \frac{\partial}{\partial z} \end{vmatrix} = \oint_{\Gamma} P\mathrm{d}x + Q\mathrm{d}y + R\mathrm{d}z. \]
利用两类曲面积分之间的联系, 还可以得到斯托克斯公式的另一形式:
\[ \iint_{\Sigma} \begin{vmatrix} \cos\alpha & \cos\beta & \cos\gamma \\ \frac{\partial}{\partial x} & \frac{\partial}{\partial y} & \frac{\partial}{\partial z} \\ P & Q & R \end{vmatrix} \mathrm{d}S = \oint_{\Gamma} P\mathrm{d}x + Q\mathrm{d}y + R\mathrm{d}z, \]
其中 \(n = (\cos\alpha, \cos\beta, \cos\gamma)\) 为有向曲面 \(\Sigma\) 在点 \((x, y, z)\) 处的单位法向量.

例 18 设 \(L\) 是柱面 \(x^2 + y^2 = 1\) 与平面 \(z = x + y\) 的交线, 从 \(z\) 轴正向往 \(z\) 轴负向看去为逆时针方向, 则曲线积分 \(\oint_L xzdx + xdy + \frac{y^2}{2}dz = \underline{\hspace{2cm}}\).

旋度 设有向量场
\[A(x, y, z) = P(x, y, z)i + Q(x, y, z)j + R(x, y, z)k.\]
若 \(P, Q, R\) 均具有一阶连续偏导数, 则向量
\[ \left(\frac{\partial R}{\partial y} - \frac{\partial Q}{\partial z}\right)i + \left(\frac{\partial P}{\partial z} - \frac{\partial R}{\partial x}\right)j + \left(\frac{\partial Q}{\partial x} - \frac{\partial P}{\partial y}\right)k \]
记作 rot A, 即
\[ \text{rot } A = \left( \frac{\partial R}{\partial y} - \frac{\partial Q}{\partial z} \right) i + \left( \frac{\partial P}{\partial z} - \frac{\partial R}{\partial x} \right) j + \left( \frac{\partial Q}{\partial x} - \frac{\partial P}{\partial y} \right) k. \]

例 19 向量场 \(A(x, y, z) = (x + y + z)i + xyj + zk\) 的旋度 \(rot A = \underline{\hspace{2cm}}\)

同步习题
见讲义第四节同步习题.